\makereporttitle{2026}{august}

\makereportpreamble

În august am continuat să lucrez pe direcțiile existente: am revizuit literatura, am reprodus diferite rezultate, am extins sau îmbunătățit rezultatele existente și am reușit să experimentez cu câteva idei noi.

Pentru studiul tranziției la haos, am reușit să înțeleg și să reproduc rezultatele obținute de John M.\ Greene în lucrarea sa din 1968, ,,\textit{Two‐Dimensional Measure‐Preserving Mappings}'', Journal of Mathematical Physics, volumul 9, nr.~5. Acesta introduce o nouă măsură asociată cu un punct fix (sau periodic) al unui sistem dinamic discret, \emph{reziduul} (,,\textit{residue}''), pe care o folosește pentru a clasifica și ierarhiza punctele periodice ale unei aplicații pe \(\reals^2\) dată de polinoame cubice. Cu toate că algoritmul folosit de el și această măsură nu sunt direct relevante pentru studiul nostru, acest articol m-a ajutat să înțeleg mai bine principiile fundamentale ale teoriei KAM.

Am reușit să reimplementez folosind JAX o parte din codurile numerice folosite pentru a studia sistemul Hénon-Heiles. Cu această ocazie am înțeles mai bine limitările metodei de integrare numerică și am identificat o metodă mai bună de a determina întreg spectrul Lyapunov pentru o traiectorie. După ce termin de implementat codurile aferente, ar trebui să pot obține o variantă finală a setului de date, pe care voi aplica algoritmul de potrivire a unei ipoteze de tip lege putere.

Pe partea de electrodinamică, am făcut mici îmbunătățiri, optimizări și reorganizări ale codului de integrare numerică pentru studiul transferului de moment cinetic orbital de la un fascicol laser Laguerre-Gauss împrăștiat de un mănunchi de electroni relativiști. Am comparat și potrivit rezultatele mele numerice cu cele obținute de un cod dezvoltat independent de către Mădălina Boca, asigurându-ne astfel că avem un studiu solid. Singurele discuții rămase sunt legate de parametri/regimul în care vrem să lucrăm, respectiv de validitatea și interpretarea formulelor pe care le folosim. Am început însă să pregătim un draft de articol pe care vrem să-l trimitem ulterior la jurnalul \textit{Matter \& Radiation at the Extremes}.

Referitor la colaborarea cu colegii de la Departamentul de Fizica Particulelor Elementare de la IFIN-HH, am mai rezolvat câteva mici probleme cu configurațiile de MadGraph pentru simularea dezintegrărilor particulei ipotetice \(S_{uu}\). Realizarea cea mai mare este pe partea de clasificare folosind metode de învățare automată. Inspirându-mă din lucrarea \href{https://arxiv.org/abs/2202.03772}{,,Particle Transformer for Jet Tagging''}, am reușit să implementez o arhitectură simplificată de tip Transformer pentru clasificarea binară de jet-uri. Rezultatele inițiale sugerează că această abordare poate obține niște numere mult mai bune, în special pe partea de eliminare a fundalului.

În colaborarea cu domnul profesor Cristian Rusu, am încercat să simplificăm și mai mult modelul teoretic pentru strategia noastră de precondiționare iterativă alternantă. Am încercat să dezvolt o metodă de identificare a rangului optim pentru precondiționarea bazată pe factorizarea Cholesky parțială. Problema se reduce la echilibrul dintre timpul inițial de calcul a matricei de precondiționare, precum și costul de aplicare per-iterație, versus numărul de iterații necesare pentru convergență. În cazul matricelor de nucleu (\textit{kernel matrix}) construite folosind nucleul Gaussian / RBF, \(\exp\left(\norm{x - y}^2/\sigma^2\right)\), avem rezultate teoretice care ne descriu spectrul. Înlocuind în formulele obținute de mine, putem să estimăm destul de precis timpul total de execuție pentru rangul curent, ceea ce înseamnă că putem identifica în timpul construcției factorizării care este rangul optim la care să ne oprim.

Luna aceasta am participat și la \href{https://sites.google.com/view/moduliandfriends2}{conferința \emph{Moduli and Friends II}} care a avut loc la \href{https://imar.ro/}{IMAR}. Am urmărit prelegeri despre topologie algebrică (cu aplicații în TQFT), teoria nodurilor, spații de moduli și alte subiecte. Totodată, am aflat că am fost acceptat la \href{https://indico.cern.ch/event/1696287/}{MLQC4FC Training School}, care se va desfășura în luna septembrie în Belgrad, Serbia. Mi-am confirmat participarea și am început demersurile pentru a solicita sprijin financiar de la UB.

\textbf{Observațiile conducătorului științific}:

\makereportpostamble{31 august 2026}
